\section*{Bab \ref{ch:b1:euclid} --- Ruang Euklides}

\begin{solution}{exo:b1:euclid:1}
$\langle u, v\rangle = 2 - 4 + 2 = 0$; $\norm u = \norm v = 3$. Adapun
vektornya \omterm{def:b1:euclid:orthogonal}{ortogonal}: sehingga sudutnya $\frac\pi2$. Cauchy--Schwarz:
$\abs 0 \leq 9$, dengan lapang.
\end{solution}

\begin{solution}{exo:b1:euclid:2}
Jabarkanlah kedua kuadrat normanya dengan kesamaan pada
\cref{thm:b1:euclid:cs} lalu tambahkanlah: maka suku silangnya saling menghapus.

Norma supnya: ambillah $x = (1, 0)$, $y = (0, 1)$. Maka $\norm{x + y}_\infty
= \norm{x - y}_\infty = 1$ dan kesamaannya akan menuntut $1 + 1 =
2(1) + 2(1) = 4$: jadi salah. Sebuah norma yang melanggar kesamaan
jajargenjang tak berasal dari \omterm{def:b1:euclid:def}{hasil kali dalam} mana pun.
\end{solution}

\begin{solution}{exo:b1:euclid:3}
$e_1 = \frac{1}{\sqrt 2}(1,1,0)$.

$w_2 = (1,0,1) - \langle (1,0,1), e_1\rangle e_1 = (1,0,1) -
\frac{1}{2}(1,1,0) = \bigl(\tfrac12, -\tfrac12, 1\bigr)$;
$\norm{w_2} = \sqrt{\tfrac32}$, sehingga $e_2 =
\frac{1}{\sqrt 6}(1, -1, 2)$.

$w_3 = (0,1,1) - \langle\cdot, e_1\rangle e_1 - \langle\cdot,
e_2\rangle e_2 = (0,1,1) - \frac12 (1,1,0) - \frac16(1,-1,2) =
\bigl(-\tfrac23, \tfrac23, \tfrac23\bigr)$; dan dengan menormalkannya, $e_3 =
\frac{1}{\sqrt 3}(-1, 1, 1)$.

(Periksa: hasil kali berpasangannya lenyap; dan masing-masingnya bernorma $1$.)
\end{solution}

\begin{solution}{exo:b1:euclid:4}
$F^\perp = \{(x,y,z) : x + y + z = 0\}$, dengan \omterm{def:b1:vspaces:free}{basis} $\bigl((1,-1,0),
(1,0,-1)\bigr)$. Dengan $e = \frac{1}{\sqrt3}(1,1,1)$: $p_F(x) =
\langle x, e\rangle e = \frac{x_1 + x_2 + x_3}{3}\,(1,1,1)$, sehingga
\[
\operatorname{Mat}(p_F) = \frac13
\begin{pmatrix} 1&1&1\\ 1&1&1\\ 1&1&1\end{pmatrix}.
\]
Untuk $x = (1,2,3)$: $p_F(x) = (2,2,2)$ dan $d(x, F) = \norm{(1,2,3)
- (2,2,2)} = \norm{(-1,0,1)} = \sqrt 2$.
\end{solution}

\begin{solution}{exo:b1:euclid:5}
Misalkan $p = \alpha(1,0,1) + \beta(0,1,1) = (\alpha, \beta, \alpha +
\beta)$. Keortogonalan $x - p$ terhadap pembangunnya:
\[
\langle x - p, (1,0,1)\rangle = (1 - \alpha) + (4 - \alpha - \beta)
= 0,
\qquad
\langle x - p, (0,1,1)\rangle = (1 - \beta) + (4 - \alpha - \beta)
= 0,
\]
yaitu $2\alpha + \beta = 5$ dan $\alpha + 2\beta = 5$: sehingga $\alpha =
\beta = \frac53$. Jadi $p_F(x) = \bigl(\frac53, \frac53,
\frac{10}{3}\bigr)$ dan
\[
d(x, F) = \norm{x - p} =
\norm{\bigl(-\tfrac23, -\tfrac23, \tfrac23\bigr)} =
\frac{2}{\sqrt 3} .
\]
Gram--Schmidt tak diperlukan karena sifat yang mendefinisikan
\omterm{def:b1:linmaps:projection}{proyeksinya} --- yaitu $x - p \perp F$ --- itu sendiri sebuah \omterm{def:b1:det:system}{sistem linear}
(``persamaan normal'') pada koefisiennya di dalam \emph{sebarang} keluarga
pembangun.
\end{solution}

\begin{solution}{exo:b1:euclid:6}
Cauchy--Schwarz dengan $g = 1$:
\[
\Bigl(\int_0^1 f \cdot 1\Bigr)^2 \leq \int_0^1 f^2 \cdot \int_0^1
1^2 = \int_0^1 f^2 ,
\]
dengan kesamaannya jika dan hanya jika $f$ sebanding dengan $1$, yaitu tetap. Adapun
versi diskretnya: di dalam $\R^n$ dengan $x = (a_1, \dots, a_n)$, $y = (1, \dots,
1)$:
$\bigl(\sum a_i\bigr)^2 \leq \norm x^2\, \norm y^2 = n \sum a_i^2$,
dengan kesamaannya jika dan hanya jika semua $a_i$ sama.
\end{solution}

\begin{solution}{exo:b1:euclid:7}
Menurut \cref{thm:b1:euclid:projection}, $E = F \oplus F^\perp$, sehingga
dimensinya bertambah: $\dim F^\perp = \dim E - \dim F$. Adapun pemuatan $F
\subseteq (F^\perp)^\perp$ seketika (karena vektor $F$
\omterm{def:b1:euclid:orthogonal}{ortogonal} terhadap segalanya di dalam $F^\perp$). Dimensinya:
$\dim(F^\perp)^\perp = \dim E - \dim F^\perp = \dim F$; dan sebuah
pemuatan dengan dimensi (hingga) yang sama merupakan kesamaan
(\cref{thm:b1:findim:subspaces}).
\end{solution}

\begin{solution}{exo:b1:euclid:8}
$A$: kolomnya bersatuan, determinannya $+1$: jadi sebuah rotasi, dengan $\cos\theta =
\frac{1}{\sqrt2}$, $\sin\theta = \frac{1}{\sqrt2}$: sehingga $\theta =
\frac\pi4$. Maka $A^8 = R_{8\pi/4} = R_{2\pi} = I$.

$B$: determinannya $\frac{1}{25}(-9 - 16) = -1$: jadi sebuah pencerminan
$S_\theta$ dengan $\cos\theta = \frac35$, $\sin\theta = \frac45$; dan
sumbunya bersudut $\frac\theta2$ dengan sumbu $x$, yaitu
garis yang berarah $\bigl(\cos\frac\theta2, \sin\frac\theta2\bigr)$
--- secara nyata sumbunya adalah
$\operatorname{Vect}\bigl((2, 1)\bigr)$, karena $B(2,1)^{\mathsf T} =
\frac15(6+4, 8-3)^{\mathsf T} = (2,1)^{\mathsf T}$. Dan sebagai pencerminan,
$B^2 = I$.
\end{solution}

\begin{solution}{exo:b1:euclid:9}
Besarannya adalah $\norm{f - (a + bX)}^2$ di dalam $C(\intcc{0}{1})$ dengan
\omterm{def:b1:euclid:def}{hasil kali dalam} \omterm{thm:b1:integration:def}{integralnya}: yang minimum tepat pada \omterm{thm:b1:euclid:projection}{proyeksi
ortogonalnya} $a + bX = p_F(f) = X - \frac16$
(\cref{ex:b1:euclid:bestapprox}). Adapun minimumnya adalah
\[
\norm{f - p_F(f)}^2 = \int_0^1 \Bigl(x^2 - x + \frac16\Bigr)^{\!2}
\dd x .
\]
Jabarkanlah: $\int_0^1 (x^2 - x + \frac16)^2 = \int_0^1 \bigl(x^4 - 2x^3
+ \frac{4}{3}x^2 - \frac13 x + \frac{1}{36}\bigr)\dd x = \frac15 -
\frac12 + \frac49 - \frac16 + \frac{1}{36} = \frac{1}{180}$.
Jadi minimumnya sama dengan $\dfrac{1}{180}$.
\end{solution}

\begin{solution}{exo:b1:euclid:10}
Misalkan $y \in \ker(u - \mathrm{id})$ (yaitu $u(y) = y$) dan $x \in E$.
Maka, dengan memakai pengawetan \omterm{def:b1:euclid:def}{hasil kali dalamnya} ($\langle u(a),
u(b)\rangle = \langle a, b\rangle$):
\[
\langle x - u(x),\, y\rangle
= \langle x, y\rangle - \langle u(x), y\rangle
= \langle x, y\rangle - \langle u(x), u(y)\rangle
= \langle x, y\rangle - \langle x, y\rangle = 0 :
\]
jadi setiap vektor $\operatorname{im}(u - \mathrm{id})$ \omterm{def:b1:euclid:orthogonal}{ortogonal}
terhadap setiap vektor yang tetap.

Karenanya $\operatorname{im}(u - \mathrm{id}) \subseteq \ker(u -
\mathrm{id})^{\perp}$, dan menurut rank--nulitas beserta
\cref{exo:b1:euclid:7}, keduanya berdimensi $\dim E - \dim\ker(u -
\mathrm{id})$: sehingga keduanya sama. Maka $E = \ker(u - \mathrm{id})
\oplus \ker(u - \mathrm{id})^\perp = \ker(u - \mathrm{id}) \oplus
\operatorname{im}(u - \mathrm{id})$.
\end{solution}

\begin{solution}{exo:b1:euclid:11}
\begin{enumerate}
  \item Andaikan $Gc = 0$ bagi sebuah kolom $c = (c_1, \dots,
        c_k)^{\mathsf T}$. Maka
        \[
        \Bigl\lVert\sum_i c_i v_i\Bigr\rVert^2
        = \sum_{i,j} c_i c_j \langle v_i, v_j\rangle
        = c^{\mathsf T} G\, c = 0,
        \]
        sehingga $\sum_i c_i v_i = 0$. Jika keluarganya \omterm{def:b1:vspaces:free}{bebas} maka ini
        memaksa $c = 0$: jadi $G$ dapat dibalik. Sebaliknya, sebuah
        hubungan yang tak sepele $\sum_j c_j v_j = 0$ memberikan, dengan mengambil
        hasil kalinya dengan setiap $v_i$, hubungan tak sepele
        $Gc = 0$: jadi $G$ singular. Maka kebebasannya $\iff \det G \neq 0$.
  \item $\langle X^{i-1}, X^{j-1}\rangle = \int_0^1 x^{i+j-2}\dd
        x = \frac{1}{i+j-1}$: sehingga matriks Gram atas $(1, X, X^2)$
        tepat merupakan \omterm{pb:b1:det:1}{matriks Hilbert} $H_3$, yang determinannya
        $\frac1{2160}$ dihitung pada soal akhir pekan
        \cref{ch:b1:det}: yang taknol, sehingga monomialnya \omterm{def:b1:vspaces:free}{bebas} ---
        seperti yang diharapkan, tetapi kini disahkan oleh sebuah bilangan.
\end{enumerate}
\end{solution}

\begin{solution}{exo:b1:euclid:12}
\begin{enumerate}
  \item $A^2 = A A = A^{\mathsf T} A = I$, sehingga $u^2 =
        \mathrm{id}$: jadi $u$ sebuah kesetangkupan, dan
        \cref{thm:b1:linmaps:projchar}~(2) memberikan $E = \ker(u -
        \mathrm{id}) \oplus \ker(u + \mathrm{id})$.
  \item Misalkan $u(x) = x$ dan $u(y) = -y$. Karena $u$ mengawetkan
        \omterm{def:b1:euclid:def}{hasil kali dalamnya},
        \[
        \langle x, y\rangle = \langle u(x), u(y)\rangle
        = \langle x, -y\rangle = -\langle x, y\rangle ,
        \]
        sehingga $\langle x, y\rangle = 0$: jadi kedua ruang eigennya
        \omterm{def:b1:euclid:orthogonal}{ortogonal}, $\ker(u + \mathrm{id}) = F^\perp$ bagi $F =
        \ker(u - \mathrm{id})$, dan $u$ merupakan kesetangkupan
        \omterm{def:b1:euclid:orthogonal}{ortogonal} terhadap $F$.
  \item $R_\theta$ setangkup jika dan hanya jika $-\sin\theta = \sin\theta$,
        yaitu $\theta \in \{0, \pi\}$: jadi \omterm{def:b1:logic:map}{pemetaan} $\pm\mathrm{id}$
        (yaitu identitas dan kesetangkupan pusat). Sedangkan setiap $S_\theta$
        setangkup: yaitu pencerminan garisnya. Inilah tepat
        kesetangkupan \omterm{def:b1:euclid:orthogonal}{ortogonal} bidangnya, dengan $F$ sama dengan
        seluruh bidangnya, $\{0\}$, atau sumbu pencerminannya.
\end{enumerate}
\end{solution}

\begin{solution}{pb:b1:euclid:1}
\textbf{1.} $\norm{t_a(x) - t_a(y)} = \norm{x - y}$; adapun \omterm{def:b1:euclid:isometry}{isometri}
\omterm{def:b1:linmaps:def}{linear} mengawetkan norma, jadi juga jarak ($\norm{g(x) - g(y)} =
\norm{g(x - y)} = \norm{x-y}$); dan komposisi \omterm{def:b1:logic:map}{pemetaan}
yang mengawetkan jarak pun mengawetkan jarak. Setiap \omterm{def:b1:euclid:isometry}{isometri} bersifat
\omterm{def:b1:logic:inj}{injektif} (karena titik yang berbeda tetap berjarak positif) dan, menurut
penggolongan di bawah, \omterm{def:b1:logic:inj}{bijektif}; sedangkan identitas dan inversnya
merupakan \omterm{def:b1:euclid:isometry}{isometri}: jadi sebuah \omterm{def:b1:structures:group}{grup}.

\textbf{2.} $\norm{f(x)} = \norm{f(x) - f(0)} = \norm{x - 0} =
\norm x$. Polarisasinya:
\[
\langle f(x), f(y)\rangle
= \frac{\norm{f(x)}^2 + \norm{f(y)}^2 - \norm{f(x) -
f(y)}^2}{2}
= \frac{\norm x^2 + \norm y^2 - \norm{x - y}^2}{2}
= \langle x, y\rangle .
\]

\textbf{3.} Dengan menjabarkannya lewat pertanyaan 2 (karena setiap hasil kali $f$
sama dengan hasil kali argumennya):
\begin{align*}
\norm{f(x+y) - f(x) - f(y)}^2
&= \norm{x+y}^2 + \norm x^2 + \norm y^2\\
&\quad - 2\langle x+y, x\rangle - 2\langle x+y, y\rangle
+ 2\langle x, y\rangle = 0 ,
\end{align*}
sebagaimana ditunjukkan periksa langsungnya; serupa pula $\norm{f(\lambda x) -
\lambda f(x)}^2 = \norm{\lambda x}^2 - 2\lambda\langle \lambda
x, x\rangle + \lambda^2\norm x^2 = 0$. Jadi $f(x + y) = f(x) +
f(y)$ dan $f(\lambda x) = \lambda f(x)$: sehingga $f$ \omterm{def:b1:linmaps:def}{linear}, dan
mengawetkan norma: jadi $f \in O(\R^2)$.

\textbf{4.} Tetapkanlah $a = f(0)$ dan $g = t_{-a}\circ f$: yaitu \omterm{def:b1:euclid:isometry}{isometri}
yang menetapkan $0$, jadi \omterm{def:b1:euclid:isometry}{isometri} \omterm{def:b1:linmaps:def}{linear} (pertanyaan 3), dan $f = t_a
\circ g$. Ketunggalannya: $t_a \circ g = t_{a'}\circ g'$ yang dinilai
di $0$ memberikan $a = a'$, lalu $g = g'$.

\textbf{5.} $(t_a\circ g)\circ(t_{a'}\circ g') = t_{a +
g(a')}\circ(g\circ g')$, karena $g\,t_{a'} = t_{g(a')}\,g$ bagi
$g$ yang \omterm{def:b1:linmaps:def}{linear}. Jadi bagian \omterm{def:b1:linmaps:def}{linear} sebuah komposisi adalah $g \circ g'$,
dan $\det(gg') = \det g \det g'$: sehingga langsung$\circ$langsung $=$
taklangsung$\circ$taklangsung $=$ langsung,
langsung$\circ$taklangsung $=$ taklangsung --- yaitu kaidah tanda $\pm1$.

\textbf{6.} $z \mapsto \eu^{\iu\theta}z$ mengirim $1$ ke
$(\cos\theta, \sin\theta)$ dan $\iu$ ke
$\iu\eu^{\iu\theta} = (-\sin\theta, \cos\theta)$: yaitu kolom
$R_\theta$. Dan $z \mapsto \eu^{\iu\theta}\conj z$ mengirim $1$ ke
$(\cos\theta, \sin\theta)$ dan $\iu$ ke $-\iu\eu^{\iu\theta} =
(\sin\theta, -\cos\theta)$: yaitu kolom $S_\theta$.

\textbf{7.} $a = 1$: maka $f = t_b$. $a \neq 1$: maka persamaan
titik tetapnya $az_0 + b = z_0$ berpenyelesaian tunggal $z_0 = b/(1 -
a)$, dan lalu
\[
f(z) - z_0 = az + b - (az_0 + b) = a(z - z_0):
\]
jadi pada kerangka yang berpusat di $z_0$, $f$ adalah perkalian dengan $a
= \eu^{\iu\arg a}$: yaitu rotasi berpusat $z_0$ dan bersudut
$\arg a$.

\textbf{8.} $f(f(z)) = a\conj{(a\conj z + b)} + b = a\conj a\,z
+ a\conj b + b = z + v$: jadi $f^2 = t_v$. Komutasinya: $f(z + v) =
a\conj z + a\conj v + b$ dan $a\conj v = a\conj{(a\conj b + b)}
= b + a\conj b = v$, sehingga $f\circ t_v = t_v\circ f$.

\emph{Kasus $v = 0$}: maka $f^2 = \mathrm{id}$. Untuk sebarang $z$,
titik tengah $m = \frac{z + f(z)}2$ memenuhi (karena \omterm{def:b1:euclid:isometry}{isometrinya}
afin, Bagian I) $f(m) = \frac{f(z) + f^2(z)}{2} = m$: jadi titik
tetapnya ada. Dengan mengonjugasikannya lewat translasi ke sebuah titik tetap,
$f$ menjadi \omterm{def:b1:euclid:isometry}{isometri} taklangsung yang \omterm{def:b1:linmaps:def}{linear}, yaitu suatu $S_\theta$
(\cref{thm:b1:euclid:plane}): jadi pencerminan pada sebuah garis.

\emph{Kasus $v \neq 0$}: maka $r = t_{-v/2}\circ f$ bersifat taklangsung dan
\[
r^2 = t_{-v/2}\,f\,t_{-v/2}\,f = t_{-v/2}\,t_{-v/2}\,f^2
= t_{-v}\,t_v = \mathrm{id}
\]
(dengan memakai komutasinya), sehingga $r$ sebuah pencerminan, dan $f =
t_{v/2}\circ r$. Adapun sumbunya sejajar $v$: karena $r$ berkomutasi dengan
$t_v$ (sebagaimana $f$ dan $t_{v/2}$), sehingga $t_v$ memetakan sumbunya (yaitu
garis tetap $r$) ke dirinya, yang memaksa $v$ mengarahkannya. Maka
$f$ merupakan \omterm{pb:b1:euclid:1}{pencerminan geser} dengan vektor geser $v/2$. Kini setiap
\omterm{def:b1:euclid:isometry}{isometri} tergolongkan: yaitu translasi atau rotasi (yang langsung),
pencerminan atau \omterm{pb:b1:euclid:1}{pencerminan geser} (yang taklangsung).

\textbf{9.} Rotasi $f(z) = az + b$, $g(z) = a'z + b'$ dengan
$a = \eu^{\iu\alpha}$, $a' = \eu^{\iu\beta}$: maka $g\circ f(z) =
a'az + (a'b + b')$ berkoefisien \omterm{def:b1:linmaps:def}{linear}
$\eu^{\iu(\alpha+\beta)}$: jadi rotasi bersudut $\alpha + \beta$
menurut pertanyaan 7, atau sebuah translasi ketika $\eu^{\iu(\alpha+\beta)} =
1$. Sedangkan pencerminan $r_i(z) = a_i\conj z + b_i$ dengan $a_i =
\eu^{\iu\theta_i}$ (yaitu sumbu bersudut $\theta_i/2$):
\[
r_2\circ r_1(z) = a_2\conj{a_1}\,z + (a_2\conj{b_1} + b_2),
\]
yang langsung dan bersudut $\theta_2 - \theta_1 = 2\bigl(\tfrac{\theta_2}2
- \tfrac{\theta_1}2\bigr)$: yaitu dua kali sudut antara sumbunya;
sedangkan sumbu yang sejajar ($\theta_1 = \theta_2$) memberikan translasi.

\textbf{10.} Rotasi berpusat $c$ dan bersudut $\theta$ merupakan
hasil kali dua pencerminan pada garis lewat $c$ yang bersudut
$\theta/2$ (yaitu pertanyaan 9, dibaca terbalik); sedangkan translasi $t_v$ adalah
hasil kali dua pencerminan pada garis sejajar yang \omterm{def:b1:euclid:orthogonal}{ortogonal} terhadap
$v$ pada jarak $\norm v/2$; sebuah pencerminan adalah satu pencerminan; dan
\omterm{pb:b1:euclid:1}{pencerminan geser} adalah pencerminan yang dikomposisikan dengan translasi,
jadi tiga. Maksimumnya: tiga.

\textbf{11.} $a = \iu$, $b = 1 + \iu$: maka $v = a\conj b + b =
\iu(1 - \iu) + 1 + \iu = (1 + \iu) + (1 + \iu) = 2 + 2\iu \neq
0$: jadi \omterm{pb:b1:euclid:1}{pencerminan geser} dengan vektor geser $v/2 = 1 + \iu$. Adapun
pencerminan $r = t_{-(1+\iu)}\circ f$ adalah $r(z) = \iu\conj z$,
yang sumbunya garis bersudut $\frac12\arg\iu = \frac\pi4$: yaitu
garis $y = x$. Jadi $f$ merupakan \omterm{pb:b1:euclid:1}{pencerminan geser} bersumbu $y =
x$ dan bervektor $(1, 1)$.

\textbf{12.} $f(z) = \iu z$ dan $g(z) = 1 + \iu(z - 1) = \iu z
+ 1 - \iu$, sehingga
\[
g\circ f(z) = \iu(\iu z) + 1 - \iu = -z + 1 - \iu :
\]
yaitu berkoefisien \omterm{def:b1:linmaps:def}{linear} $-1 = \eu^{\iu\pi}$, jadi rotasi bersudut
$\pi$ (yakni setengah putaran), berpusat $z_0 = \frac{1 - \iu}{1 - (-1)}
= \frac{1 - \iu}{2}$.

\textbf{13.} $r_2\circ r_1(z) = \iu\,\conj{\conj z} = \iu z$:
yaitu rotasi bersudut $\frac\pi2$ terhadap $0$. Adapun sumbunya (yakni sumbu
real, bersudut $0$; dan garis $y = x$, bersudut $\frac\pi4$) bertemu pada
sudut $\frac\pi4$, dan dua kali itu adalah $\frac\pi2$: jadi pertanyaan 9
tertegaskan.

\textbf{14.} Dengan $f = t_a\circ g$, $g$ yang \omterm{def:b1:linmaps:def}{linear}, dan $\sum_i
\lambda_i = 1$:
\[
f\Bigl(\sum_i \lambda_i x_i\Bigr)
= a + \sum_i \lambda_i\,g(x_i)
= \sum_i \lambda_i\bigl(a + g(x_i)\bigr)
= \sum_i \lambda_i f(x_i) .
\]

\textbf{15.} Untuk $h \in G$, dengan memakai pertanyaan 14 (karena koefisien
$\frac1{\abs G}$ berjumlah $1$):
\[
h(c) = \frac{1}{\abs G}\sum_{g \in G} h\bigl(g(x_0)\bigr)
= \frac{1}{\abs G}\sum_{g' \in G} g'(x_0) = c ,
\]
karena $g \mapsto hg$ merupakan bijeksi $G$ pada dirinya sendiri.

\textbf{16.} \omterm{def:b1:complex:field}{Konjugat} $t_{-c}\,g\,t_{c}$ ($g \in G$) membentuk
sebuah \omterm{def:b1:structures:group}{grup} \omterm{def:b1:euclid:isometry}{isometri} yang menetapkan $0$, jadi \omterm{def:b1:euclid:isometry}{isometri} yang \emph{\omterm{def:b1:linmaps:def}{linear}}
(pertanyaan 3): yaitu \omterm{def:b1:structures:subgroup}{subgrup} hingga pada $O(\R^2)$. Jika
suatu $s \in G$ bersifat taklangsung, maka \omterm{def:b1:logic:map}{pemetaan} $g \mapsto sg$ mengirim $G^+$
secara \omterm{def:b1:logic:inj}{injektif} ke dalam $G\setminus G^+$ dan $h \mapsto s^{-1}h$ merupakan
inversnya (lewat kaidah tanda pertanyaan 5): sehingga $\abs G = 2\,\abs{G^+}$;
kalau tidak $G = G^+$.

\textbf{17.} Jika $G^{+} = \{\mathrm{id}\}$, maka ia $C_1$.
Kalau tidak tulislah unsurnya sebagai $R_\theta$, $\theta \in
\intco{0}{2\pi}$, lalu misalkan $\theta_0$ sudut positif
terkecil yang muncul. Untuk $R_\theta \in G^{+}$, pembagian Euclid
$\theta = k\theta_0 + \rho$ dengan $0 \leq \rho < \theta_0$ memberikan
$R_\rho = R_\theta R_{\theta_0}^{-k} \in G^{+}$, sehingga $\rho = 0$
lewat keminimalannya: jadi $G^{+} = \langle R_{\theta_0}\rangle$. Adapun \omterm{def:b1:arith:divides}{membagi}
$2\pi$ oleh $\theta_0$ dengan cara yang sama menunjukkan $2\pi = n\theta_0$:
sehingga $G^{+} \cong C_n$, yang dibangun oleh rotasi bersudut
$\frac{2\pi}n$.

\textbf{18.} $G = G^{+} \cup sG^{+}$ menurut pencacahan
pertanyaan 16. Dengan $s(z) = a\conj z$ dan $r(z) = \omega z$
(yaitu \omterm{def:b1:linmaps:forms}{bentuk linear}, setelah pertanyaan 16; $\abs a = \abs\omega = 1$):
\[
s r s(z) = a\,\conj{\omega\,a\conj z}
= a\conj\omega\conj a\,z = \conj\omega\,z = r^{-1}(z) .
\]
Setiap unsur $sr^k$ pada $sG^{+}$ bersifat taklangsung dan memenuhi
$(sr^k)^2 = (sr^ks)\,r^k = r^{-k}r^k = \mathrm{id}$: yaitu
\omterm{def:b1:euclid:isometry}{isometri} taklangsung yang involutif dan menetapkan $c$, jadi sebuah pencerminan. Sehingga
$G$ terdiri atas $n$ rotasi dan $n$ pencerminan dengan
hubungan $r^n = s^2 = \mathrm{id}$, $srs = r^{-1}$: yaitu
\omterm{pb:b1:euclid:1}{grup dihedral} $D_n$.

\textbf{19.} Menggabungkannya: sebuah \omterm{def:b1:structures:group}{grup} hingga berisi \omterm{def:b1:euclid:isometry}{isometri} bidang
menetapkan sebuah titik $c$ (pertanyaan 15), tereduksi menjadi \omterm{def:b1:structures:subgroup}{subgrup} hingga
pada $O(2)$ (pertanyaan 16), dan berupa $C_n$ jika ia memuat hanya
rotasi (pertanyaan 17), dan $D_n$ kalau tidak (pertanyaan 18):
yaitu \emph{teorema Leonardo}.

\textbf{20.} Sebuah translasi $t_v \neq \mathrm{id}$ berpangkat
$t_{kv}$, yang semuanya berbeda ($kv$ berbeda berpasangan untuk $v \neq 0$):
jadi berorde takhingga. Sedangkan \omterm{pb:b1:euclid:1}{pencerminan geser} $f$ mempunyai $f^2 = t_v$ dengan $v
\neq 0$: berorde takhingga lagi. Tak satu pun dari keduanya muat di dalam \omterm{def:b1:structures:group}{grup} hingga
--- yang selaras dengan pertanyaan 15--19, yang hanya menghasilkan
rotasi dan pencerminan.

\textbf{21.} Kesemua $2n$ \omterm{def:b1:logic:map}{pemetaannya} mengawetkan \omterm{def:b1:logic:sets}{himpunan} titik sudutnya: $\omega^k
\omega^j = \omega^{k+j}$ dan $\omega^k\conj{\omega^j} =
\omega^{k-j}$; dan karena \omterm{def:b1:euclid:isometry}{isometri}, semuanya mengawetkan segi banyaknya
(karena ruas garisnya pergi ke ruas garis). Sebaliknya, sebuah kesetangkupan mengawetkan
barisentrum $0$ titik sudutnya (pertanyaan 14), jadi bersifat \omterm{def:b1:linmaps:def}{linear},
dan mempermutasikan titik sudutnya (karena itulah titik segi banyaknya
yang berjarak maksimal dari $0$). Adapun \omterm{def:b1:euclid:isometry}{isometri} \omterm{def:b1:linmaps:def}{linear} yang mengirim
titik sudut $1$ ke $\omega^k$ adalah $z \mapsto \omega^k z$ bila langsung,
dan $z \mapsto \omega^k\conj z$ bila taklangsung (karena matriksnya
ditentukan oleh satu kolom dan tandanya): jadi paling banyak $2n$
kesetangkupan, sehingga tepat $D_n$ di atas.

\textbf{22.} $C_1$: yaitu segitiga sembarang (tanpa kesetangkupan
yang tak sepele). $D_1$: yaitu segitiga samakaki yang bukan samasisi (dengan satu
pencerminan). $D_2$: yaitu persegi panjang yang bukan persegi (dengan identitas,
setengah putaran terhadap pusatnya, dan kedua pencerminan sumbunya). $C_3$: yaitu
triskelion --- tiga lengan bengkok yang sebangun dan terpasang pada \omterm{prop:b1:reals:intervals}{selang}
$120$ derajat; karena bengkoknya membunuh setiap pencerminan tetapi menyisakan
rotasi berorde $3$.

\textbf{23.} Dengan $\omega = \iu$: yaitu rotasi $z \mapsto z$,
$\iu z$, $-z$, $-\iu z$ (bersudut $0, \frac\pi2, \pi,
\frac{3\pi}2$), dan pencerminan
\[
z \mapsto \conj z \ (\text{sumbu } y = 0), \quad
\iu\conj z \ (y = x), \quad
-\conj z \ (x = 0), \quad
-\iu\conj z \ (y = -x):
\]
yaitu kedua diagonal dan kedua garis tengah perseginya --- jadi
\omterm{pb:b1:euclid:1}{grup dihedral} $D_4$, berorde $8$.

\textbf{24.} Tulislah $a = \eu^{\iu\theta} \neq 1$: maka $f(z) = az +
c_1(1 - a)$ dan $g(z) = az + c_2(1 - a)$. Lalu
\[
f\circ g(z) - g\circ f(z)
= (1-a)\bigl(ac_2 + c_1 - ac_1 - c_2\bigr)
= -(1 - a)^2(c_2 - c_1) \neq 0 ,
\]
sehingga $f g \neq g f$. Kedua komposisinya berkoefisien \omterm{def:b1:linmaps:def}{linear}
$a^2$, sehingga $(fg)\circ(gf)^{-1}$ berkoefisien \omterm{def:b1:linmaps:def}{linear} $1$: jadi ia
translasi sebesar $fg(z) - gf(z)$, yaitu tetapan taknol di atas.
Adapun \omterm{def:b1:structures:group}{grup} yang memuat $f$ dan $g$ memuat translasi ini dan
semua pangkatnya: jadi ia takhingga. Dua pusat rotasi sudah satu
kelebihan bagi \omterm{def:b1:structures:group}{grup} hingga --- yaitu jantung geometri teorema
Leonardo.

\textbf{25.} (i) Pertanyaan 3--4 menunjukkan bahwa setiap \omterm{def:b1:logic:map}{pemetaan}
yang mengawetkan jarak bersifat afin dengan bagian \omterm{def:b1:linmaps:def}{linear} yang \omterm{def:b1:euclid:orthogonal}{ortogonal}, dan pertanyaan 15 memakukan
setiap \omterm{def:b1:structures:group}{grup} hingga pada sebuah titik tetap: sehingga kedua hasilnya mengubah
geometri metriknya menjadi aljabar \omterm{def:b1:linmaps:def}{linear} di titik asalnya. (ii) Adapun
daftar lengkapnya adalah translasi, rotasi, pencerminan, dan \omterm{pb:b1:euclid:1}{pencerminan
geser}; dengan determinan bagian \omterm{def:b1:linmaps:def}{linearnya} memisahkan yang langsung
dari yang taklangsung, sedangkan keberadaan titik tetapnya memisahkan
kedua jenisnya di dalam tiap keparian. (iii) Menurut pertanyaan 9, dua
pencerminan berkomposisi menjadi sebarang rotasi atau translasi, sehingga
pencerminannya membangun seluruh \omterm{def:b1:structures:group}{grupnya} --- dan paling banyak tiga sudah cukup
bagi sebarang \omterm{def:b1:euclid:isometry}{isometri}. (iv) Pada skala hingganya hanya dua keluarga
yang bertahan, yaitu \omterm{def:b1:structures:group}{grup} siklik dan dihedral, dan itulah sebabnya hiasan
rozet hadir tepat dalam dua macam (dengan atau tanpa sumbu
cermin). Bagian II membuktikan \emph{penggolongan \omterm{def:b1:euclid:isometry}{isometri}
bidang}; sedangkan Bagian IV membuktikan \emph{teorema Leonardo}.

\end{solution}
